%% BEAMER THEME ITHACA 2013: Main tex file for compiling
%$ Compile this file. 
%%
%% Copyright 2013 by Raul Gomez
%% This file may be distributed and/or modified
%% 	1. under the LaTeX Project Public License and/or
%% 	2. under the GNU Public License.
%% 
%% This theme is based on Flip's Theme by Flip Tanedo.

%% Please see notes.txt for comments on Beamer Theme Ithaca 2013
%% By default, this template is meant to be run with XeLaTeX (for fonts)
%% To run in PDFLaTeX, remove fontspec and any font commands

%% Discussion of Beamer vs XeLaTeX vs LuaLaTeX
%% http://tex.stackexchange.com/questions/29497/xelatex-preventing-beamer-from-using-different-backgrounds



%\documentclass[12 pt,hyperref={pdfpagelabels=false}]{beamer}
%\documentclass[aspectratio=169]{beamer}



\usetheme[
%	bullet=square,           % Default option: circle
%	nobigpagenumber,         % No big circled page number on lower right
%       nofootline,		 % No footline
	grey                    % For a neutral grey background color
%	watermark=BG_both,	 % png file for the watermark
	]{Ithaca}



%%%%%%%%%%
% FONTS %
%%%%%%%%%%

%% Default font: lmodern, doesn't require fontspec % solves some default warnings
\usepackage[T1]{fontenc}
\usepackage{lmodern}	
\usepackage{mathptmx}
\usepackage{frutiger}		

%% Protects fonts from Beamer screwing with them
%% http://tex.stackexchange.com/questions/10488/force-computer-modern-in-math-mode
\usefonttheme{professionalfonts}



%%%%%%%%%%%%%%%%%%%%%%%%
% Usual LaTeX Packages %
%%%%%%%%%%%%%%%%%%%%%%%%

\usepackage{amsmath}
\usepackage{amsfonts}
\usepackage{amssymb}
\usepackage{graphicx}
\usepackage{mathrsfs} 			% For Weinberg-esque letters
\usepackage{cancel}			% For "SUSY-breaking" symbol
\usepackage{slashed}   		        % for slashed characters in math mode
\usepackage{bbm}                	% for \mathbbm{1} (unit matrix)
\usepackage{amsthm}			% For theorem environment
\usepackage{multirow}			% For multi row cells in table
\usepackage{young}			% For Young Tableaux
\usepackage{xspace}			% For spacing after commands
\usepackage{wrapfig}			% for Text wrap around figures
\usepackage{framed}
\usepackage{pgf}
\usepackage[table]{xcolor}
\usepackage{tikz}
\usepackage{hyperref}
\usepackage{float}



\graphicspath{{images/}}		% Put all images in this directory. Avoids clutter.

\newcommand{\cG}{\mathcal{T}}
\newcommand{\supp}{\operatorname{supp}}
\newcommand{\Span}{\operatorname{Span}}
\newcommand{\R}{\mathbb{R}}
\newcommand{\C}{\mathbb{C}}
\newcommand{\N}{\mathbb{N}}
\newcommand{\Z}{\mathbb{Z}}
\newcommand{\RP}{\mathbb{RP}}
\newcommand{\g}{\mathfrak{g}}
\newcommand{\n}{\mathfrak{n}}
\newcommand{\m}{\mathfrak{m}}
\newcommand{\p}{\mathfrak{p}}
\renewcommand{\sl}{\mathfrak{sl}}
\renewcommand{\S}{\mathscr{S}}
\newcommand{\Hom}{\operatorname{Hom}}
\newcommand{\SL}{\operatorname{SL}}
\newcommand{\PGL}{\operatorname{PGL}}
\renewcommand{\a}{\mathfrak{a}}
\newcommand{\ind}{\operatorname{ind}}
\newcommand{\Ext}{\operatorname{Ext}}
\newcommand{\Tor}{\operatorname{Tor}}
\newcommand{\hot}{\widehat{\otimes}}
\newcommand{\Fre}{{Fr\'{e}chet \,}}
\newcommand{\Sc}{{\mathcal S}}
\newcommand{\sgn}{\operatorname{sgn}}
\newcommand{\Orb}{\mathcal{O}}
\newcommand{\Cov}{\operatorname{Cov}}
\newcommand{\Var}{\operatorname{Var}}
\newcommand{\Tr}{\operatorname{Tr}}
\renewcommand{\N}{\mathcal{N}}
\newcommand{\diag}{\operatorname{diag}}
\newcommand{\E}{\mathbb{E}}


%%%%%%%%%%%%%%%%%%%%%%%%%%%%%%%%%%%%%%%%%%%%%%
% Theorem, Definition and Remark Environments%
%%%%%%%%%%%%%%%%%%%%%%%%%%%%%%%%%%%%%%%%%%%%%%

\theoremstyle{plain}
%\newtheorem{proposition}[theorem]{Proposition}
%\newtheorem{conjecture}[theorem]{Conjecture}
%\newtheorem{claim}[theorem]{Claim}
%\newtheorem{algorithm}[theorem]{Algorithm}
%\newtheorem{axiom}[theorem]{Axiom}
%\newtheorem{criterion}[theorem]{Criterion}

\theoremstyle{definition}
%\newtheorem{condition}[theorem]{Condition}
%\newtheorem{exercise}[theorem]{Exercise}
%\newtheorem{notation}[theorem]{Notation}
%\newtheorem{case}[theorem]{Case}
%\newtheorem{question}[theorem]{Question}

\theoremstyle{remark}
%\newtheorem{remark}[theorem]{Remark}
%\newtheorem{remarks}[theorem]{Remarks}
%\newtheorem{summary}[theorem]{Summary}
%\newtheorem{observation}[theorem]{Observation}
%\newtheorem{conclusion}[theorem]{Conclusion}
%\newtheorem{acknowledgement}[theorem]{Acknowledgement}


\newcounter{saveenumi}
\newcommand{\seti}{\setcounter{saveenumi}{\value{enumi}}}
\newcommand{\conti}{\setcounter{enumi}{\value{saveenumi}}}
%%%%%%%%%%%%%%%%%%%%%%%%%%%%%%%%%%%%%%
% Section and subsection adjustments %
%%%%%%%%%%%%%%%%%%%%%%%%%%%%%%%%%%%%%%
\AtBeginSection{}
\AtBeginSubsection{}	 




%%%%%%%%%%%%
% Preamble %
%%%%%%%%%%%%

\title[MA2003B]{Regresión lineal múltiple}
\subtitle{MA2003B Aplicación de métodos multivariados \\ en ciencia de datos}
\institute[Tec de Monterrey]{EIC}
\date{Campus Monterrey}
\logo{\includegraphics[height=1.0cm]{tecnologico-de-monterrey-blue}}
